\documentclass{article}

\title{Comparing Specialty Portrait Lenses}
\subtitle{(Draft 2026-10-01)
}
\author{LAK132}

\begin{document}
\begin{lstlisting}[language=page-meta]
{["description"]="STF/DS/APD vs SF vs DC vs VFC vs anamorphic",
 ["image"]="resources/DSC01570_edited.jpg"}
\end{lstlisting}

\maketitle

Have you ever wondered what the differences between
\abbr{STF}{Smooth Trans Focus (Apodized, Minolta/Sony)}/\abbr{DS}{Defocus Smoothing (Apodized, Canon)}/\abbr{APD}{Apodized (Fuji)},
\abbr{SF}{Soft Focus},
\abbr{DC}{Defocus Control (Nikon)},
\abbr{VFC}{Variable Field Curvature (Minolta)},
and anamorphic lenses are?

\begin{figure}
	\caption{Sony a7SII, Sony A 135mm f/2.8 T4.5 STF, ISO1000 1/160s T4.5}
	\includegraphics{\postpath resources/DSC01570_edited.jpg}
\end{figure}


\section{Standard Optics}

\charq[LAK132]{\webroot resources/pfp/lak132.png}{
What even counts as a portrait lens anyway?

Let's go with ``prime'' and ``close to 135mm''
(because 135mm is \textit{obviously} the best focal length)
}

\subsection{Lenses}

A few from my collection:

\begin{itemize}
	\item Pentax Super-Multi-Coated TAKUMAR 135mm f/2.5 6-element
	\item Sony A Carl Zeiss Sonnar T* 135mm f/1.8 ZA
	\item Zeiss Tele-Tessar 135/4 (Zeiss Ikon Contaflex 126)
\end{itemize}

\charq[LAK132]{\webroot resources/pfp/lak132.png}{
This page probably wouldn't load if I attempted to list every single standard
``portrait'' lens ever made lmao
}

\subsection{Testing}

\subsubsection{Pentax Super-Multi-Coated TAKUMAR 135mm f/2.5 6-element}

\begin{figure}
	\includegraphics[width=27em]{\postpath resources/DSC00591.JPG}
	\includegraphics[width=27em]{\postpath resources/DSC00597.JPG}
	\caption{Sony a7RV, Pentax Super-Multi-Coated TAKUMAR 135mm f/2.5 6-element}
\end{figure}

\subsubsection{Sony A Carl Zeiss Sonnar T* 135mm f/1.8 ZA}

\section{Apodization}

\subsection{Lenses}

Very few lenses with apodization filters have ever been made.

\begin{itemize}
	\item Sony/Minolta A 135mm f/2.8 T4.5 \textbf{\abbr{STF}{Smooth Trans Focus}}
	\item Sony FE 100mm f/2.8 T5.6 \textbf{\abbr{STF}{Smooth Trans Focus}} GM OSS
	\item LAOWA 105mm f/2 T3.2 \textbf{\abbr{STF}{Smooth Trans Focus}}
	\item Fujinon Super EBC XF 56mm f/1.2 T1.7 R \textbf{\abbr{APD}{Apodized}}
	\item Canon RF 85mm f/1.2 T2 L USM \textbf{\abbr{DS}{Defocus Smoothing}}
\end{itemize}

Something you may immediately notice about the names of these lenses is that
they list both an f-number \textit{and} a T-number.
For example, the Sony FE 100mm f/\textbf{2.8} T\textbf{5.6} STF has an
f-number of 2.8 and a T-number of 5.6.
Usually a lens's f-number (focal length divided by diameter of the entrance
pupil (\( N = f / D \)), often used as a measure of a lens's depth of field)
and T-number (measure of light transmission) are close enough in value that
lens manufacturers only report one or the other (f-number for photographic
lenses, T-number for cine lenses).
In the case of apodized lenses however, these values can and do differ quite
significantly.

\subsection{Testing}

\subsubsection{Sony A 135mm f/2.8 T4.5 STF}

\begin{figure}
	\includegraphics[width=27em]{\postpath resources/DSC01464.JPG}
	\includegraphics[width=27em]{\postpath resources/DSC01476.JPG}
	\caption{Sony a7RV, Sony A 135mm f/2.8 T4.5 STF}
\end{figure}

\section{Soft Focus}

Soft focus lenses typically fall into 1 of 2 categories:
\begin{itemize}
	\item Lenses that use optical formulation tricks to prevent sharp focus of the image, and
	\item Lenses that use aperture tricks to prevent sharp focus of the image.
\end{itemize}
Lenses with adjustable soft focus will usually be the former,
while the latter will generally require the lens to be stopped down quite
considerably to prevent the soft focus effect.

\subsection{Lenses}

Many many many different soft focus lenses exist. A few notable ones:

\begin{itemize}
	\item Canon EF 135mm f/2.8 \textbf{Soft Focus} (optical soft focus)
	\item Minolta AF 100mm f/2.8 \textbf{Soft Focus} (optical soft focus)
	\item Fujinon EBC M42 \textbf{\abbr{SF}{Soft Focus}} 85mm f/4 (aperture soft focus)
\end{itemize}

\subsection{Testing}

\subsubsection{Canon EF 135mm f/2.8 Soft Focus}

\subsubsection{Minolta AF 100mm f/2.8 Soft Focus}

\subsubsection{Fujinon EBC M42 SF 85mm f/4}

The trypophobia lens!

\section{Defocus Control}

\subsection{Lenses}

The Nikon DC-NIKKORs seems to be the only examples of this type of lens.

\begin{itemize}
	\item Nikon F AF \textbf{\abbr{DC}{Defocus Control}}-NIKKOR 135mm f/2
	\item Nikon F AF \textbf{\abbr{DC}{Defocus Control}}-NIKKOR 135mm f/2 D
	\item Nikon F AF \textbf{\abbr{DC}{Defocus Control}}-NIKKOR 105mm f/2 D
\end{itemize}

\subsection{Testing}

\subsubsection{Nikon F AF DC-NIKKOR 105mm f/2 D}

\begin{figure}
	\includegraphics[width=18em]{\postpath resources/DSC_0338.JPG}
	\includegraphics[width=18em]{\postpath resources/DSC_0339.JPG}
	\includegraphics[width=18em]{\postpath resources/DSC_0340.JPG}
	\caption{Nikon D200, Nikon F AF DC-NIKKOR 105mm f/2 D (157mm full frame equivalent)}
\end{figure}

Unfortunately my particular copy of this lens appears to have some issues with
chromatic aberrations, that other people who have owned this lens before have
told me were not present on their copies.

\subsubsection{Nikon F AF DC-NIKKOR 135mm f/2 D}

\section{Variable Field Curvature}

\subsection{Lenses}

\begin{itemize}
	\item Minolta MC \textbf{\abbr{VFC}{Variable Field Curvature}} Rokkor 24mm f/2.8
	\item Minolta MD \textbf{\abbr{VFC}{Variable Field Curvature}} Rokkor-X 24mm f/2.8
\end{itemize}

Considering the relatively wide focal length of 24mm of these two lenses,
they're probably not great for portraiture.
Nevertheless, their rather unique focus control method felt relevant,
so I'm including them anyway.
\charq[LAK132]{\webroot resources/pfp/lak132.png}{
A 24mm lens is a portrait lens if the portrait calls for a 24mm lens
}

\subsection{Testing}

\subsubsection{Minolta MD VFC Rokkor-X 24mm f/2.8}

\section{Anamorphic}

\charq[Kicodora]{\webroot resources/pfp/kicodora2.png}{
Really?

Anamorphic?

The lenses pretty much exclusively used for wide screen movies?
}
\charq[LAK132]{\webroot resources/pfp/lak132.png}{
ANY LENS IS A PORTRAIT LENS IF THE PORTRAIT CALLS FOR IT
}

In addition to the regular focal length and f-stop/T-stop numbers included
on other lenses, anamorphics also include a compression factor,
typically something like 1.33x, 1.6x, 1.8x, etc.
This number represents how much the image is compressed horizontally,
and how much the image must be stretched back out to make it look ``normal''
again.
This horizontal compression allows a wide screen image to be recorded on an
otherwise relatively narrow medium, without wasting any of the medium's
resolution, as you would otherwise have to do if you cropped the image to the
same aspect ratio.

The major stylistic side effects of this compression process is ``bokeh balls''
(circles of confusion created by out of focus specular highlights) end up
as ovals instead of circles, and the compression optics have a tendency to
create horizontal flares that cover most, if not all,
of the width of the image.

\subsection{Lenses}

This is another lens category with far too many lenses to mention here,
so we'll limit the list to just the full frame 135mm lenses that I can easily
find info about:

\begin{itemize}
	\item Sirui Venus 135mm T2.9 \textbf{1.8x}
	\item Sirui IronStar 135mm T1.9 \textbf{1.5x}
	\item Viltrox EPIC 135mm T2.4 \textbf{1.33x}
	\item Blazar Mantis 135mm T3.2 \textbf{1.33x}
	\item Cooke Anamorphic/i FF 135mm T2.3 \textbf{1.8x}
\end{itemize}

\subsection{Testing}

\subsubsection{Sirui Venus 135mm T2.9 1.8x}

\section{The 135mm shootout}

\begin{itemize}
	\item Sony A Carl Zeiss Sonnar T* 135mm f/1.8 ZA
	\item Sony A 135mm f/2.8 T4.5 \textbf{\abbr{STF}{Smooth Trans Focus}}
	\item Nikon F AF \textbf{\abbr{DC}{Defocus Control}}-NIKKOR 135mm f/2 D
	\item Canon EF 135mm f/2.8 \textbf{Soft Focus}
	\item Sirui Venus 135mm T2.9 \textbf{1.8x}
\end{itemize}

\end{document}
